\documentclass{article}

% Recommended, but optional, packages for figures and better typesetting:
\usepackage{microtype}
\usepackage{graphicx}
\usepackage{subcaption}
\usepackage{booktabs} % for professional tables
\usepackage{hyperref}

\newcommand{\theHalgorithm}{\arabic{algorithm}}

% Use the following line for the initial blind version submitted for review:
\usepackage{icml2026}

\usepackage{amsmath}
\usepackage{amssymb}
\usepackage{mathtools}
\usepackage{amsthm}

% if you use cleveref..
\usepackage[capitalize,noabbrev]{cleveref}

%%%%%%%%%%%%%%%%%%%%%%%%%%%%%%%%
% THEOREMS
%%%%%%%%%%%%%%%%%%%%%%%%%%%%%%%%
\theoremstyle{plain}
\newtheorem{theorem}{Theorem}[section]
\newtheorem{proposition}[theorem]{Proposition}
\newtheorem{lemma}[theorem]{Lemma}
\newtheorem{corollary}[theorem]{Corollary}
\theoremstyle{definition}
\newtheorem{definition}[theorem]{Definition}
\newtheorem{assumption}[theorem]{Assumption}
\theoremstyle{remark}
\newtheorem{remark}[theorem]{Remark}

\icmltitlerunning{SO-LoRA: Parameter-Efficient Multi-Task Model Merging}

\begin{document}

\twocolumn[
  \icmltitle{SO-LoRA: Parameter-Efficient Multi-Task Model Merging via \\
    Subspace-Disjoint Orthogonal Low-Rank Adaptation}

  \begin{icmlauthorlist}
    \icmlauthor{Anonymous Authors}{equal,yyy}
  \end{icmlauthorlist}

  \icmlaffiliation{yyy}{Department of Computer Science, University of Machine Learning, Location, Country}
  \icmlcorrespondingauthor{Anonymous Authors}{contact@anonymous.edu}

  \icmlkeywords{Model Merging, PEFT, LoRA, SAM, Flat Minima, Weight Space Orthogonality}

  \vskip 0.3in
]

\printAffiliationsAndNotice{}

\begin{abstract}
  Parameter-efficient model merging has emerged as a promising, training-free paradigm for combining diverse capabilities from task-specific expert models. However, when merging independently trained Low-Rank Adaptation (LoRA) adapters via parameter averaging or addition, task interference in both the weight and activation spaces leads to severe representation collapse. In this work, we propose \textbf{Subspace-Disjoint Orthogonal LoRA (SO-LoRA)}, a novel parameter-efficient fine-tuning and model-merging framework that mathematically guarantees zero interference between task-specific update directions. Specifically, we deterministically pre-partition the low-rank output projection space using a shared orthogonal basis, restricting each task's low-rank updates to mutually orthogonal, disjoint subspaces. We couple this structural orthogonality with \textbf{Sharpness-Aware Minimization (SAM)} during independent training, ensuring that each adapter lands in wide, robust flat minima. This dual geometric protection---subspace-disjoint weight orthogonality and loss landscape flatness---enables lossless model merging. We evaluate SO-LoRA on Split CIFAR-10 classification, achieving up to 30\% higher merged multi-task accuracy over standard LoRA merging, and significantly outperforming SAM-only merging. Our approach offers a simple, highly scalable, and communication-free solution for collaborative modular model customization.
\end{abstract}

\section{Introduction}
\label{intro}

The pretrain-finetune paradigm has revolutionized deep learning, allowing models to acquire specialized knowledge across a wide array of downstream tasks \cite{devlin2018bert, vaswani2017attention}. To mitigate the prohibitive computational and storage costs of full-parameter fine-tuning, Parameter-Efficient Fine-Tuning (PEFT) methods, particularly Low-Rank Adaptation (LoRA) \cite{hu2021lora}, have been widely adopted. LoRA represents weight updates as a product of low-rank matrices, dramatically reducing the number of learnable parameters while maintaining competitive performance.

As decentralized and localized training gains prominence, there is an increasing demand to compose independently trained task-specific models into a single unified network without incurring retraining costs. This has motivated the study of \textit{model merging} \cite{wortsman2022model, li2023deep}, which aims to fuse multiple neural network weights directly. The simplest and most popular merging approach is Task Arithmetic \cite{ilharco2022editing}, which directly averages or adds the task-specific parameter updates to the base model.

Despite its simplicity, direct parameter averaging often results in catastrophic representation collapse \cite{yadav2023ties}. This is due to \textit{task interference}: independently trained networks develop incompatible coordinate systems and gradient directions, causing destructive interference when their weights are summed or averaged. This issue is particularly acute in LoRA merging, where the low-rank projection directions of different adapters are highly misaligned and easily overwrite each other's activation spaces.

Prior works have attempted to address this challenge by aligning permutation symmetries post-hoc \cite{ainsworth2022git, mi2022permutation}, pruning parameter updates \cite{yadav2023ties, yu2024dare}, or performing expensive test-time synergistic adaptation \cite{symerge2025steering}. Others have shown that optimizing the training phase using Sharpness-Aware Minimization (SAM) \cite{foret2021sam} can significantly enhance model mergeability by steering task-specific parameters into wide, robust flat minima of the loss landscape \cite{satalr2026flatter, saortho2026saortho}. However, while flatness provides local robustness, it does not prevent global structural overlap between different models' representational subspaces.

To resolve this fundamental limitation, we present \textbf{Subspace-Disjoint Orthogonal LoRA (SO-LoRA)}. Instead of addressing interference post-hoc, we proactively design the fine-tuning process to guarantee zero interference between task-specific updates. 
We generate a shared, deterministic orthogonal basis using a fixed random seed. For $K$ tasks and a rank-$r$ LoRA adaptation, we pre-partition this basis into $K$ disjoint, mutually orthogonal subspaces. During fine-tuning, each task's low-rank projection is strictly constrained to its designated subspace. Since the columns of the low-rank output projection matrices of any two distinct tasks are mutually orthogonal by design, their updates reside in disjoint representational directions, ensuring perfect linear addition with zero interference.

Furthermore, we couple this structural orthogonality with SAM during training. This double geometric protection---loss landscape flatness and weight-space subspace disjointness---guarantees that when the expert models are merged, they combine seamlessly without destroying each other's decision boundaries.

We evaluate SO-LoRA on Split CIFAR-10, where a pre-trained ResNet-18 model is fine-tuned on disjoint task subsets (classes 0--4 and 5--9). Our experiments demonstrate that:
\begin{itemize}
  \item Standard LoRA merging suffers from significant accuracy drop due to weight-space interference.
  \item Merging on orthogonal subspaces (SO-LoRA) dramatically boosts performance, allowing the merged model to maintain high multi-task accuracy.
  \item Combining SO-LoRA with SAM achieves the highest overall accuracy, demonstrating the complementary power of weight-space orthogonality and loss landscape flatness.
  \item Our approach is entirely decentralized, communication-free, and does not require expensive test-time optimization or parameter tuning.
\end{itemize}

\section{Related Work}
\label{related}

\textbf{Model Merging and Task Arithmetic:} Directly averaging weights or adding parameter updates (Task Arithmetic) has emerged as an attractive, training-free path to compose multiple capabilities \cite{wortsman2022model, ilharco2022editing}. However, direct addition is prone to destructive interference because independent fine-tunings produce divergent parameter magnitudes and conflicting sign agreements. Methods like TIES-Merging \cite{yadav2023ties} and DARE \cite{yu2024dare} attempt to resolve this post-hoc by pruning redundant parameters, dropping low-magnitude updates, and resolving sign conflicts. Git Re-Basin \cite{ainsworth2022git} aligns model weight permutations using coordinate-space symmetries. While effective, these post-hoc approaches can be computationally complex and often fail under extreme task divergence.

\textbf{Merging Low-Rank Adapters:} Merging task-specific LoRA adapters has attracted significant attention. ZipLoRA \cite{po2023orthogonal} uses a soft-regularization penalty during joint training to promote low-rank update orthogonality. Core Space Merging \cite{po2023orthogonal} and KnOTS seek to reduce representational overlap post-hoc. Unlike these methods, SO-LoRA pre-partitions the low-rank representational space prior to training, requiring zero communication or joint training, which makes it ideal for decentralized settings.

\textbf{Flat Minima and SAM in Model Merging:} Sharpness-Aware Minimization (SAM) \cite{foret2021sam} trains models to lie in flat, low-curvature regions of the loss landscape. Flat minima are robust to parameter perturbations, which directly improves their mergeability. Recent works have shown that SAM training during fine-tuning dramatically enhances LoRA model merging (SATA-LR) \cite{satalr2026flatter} and full-weight orthogonal merging (SA-Ortho) \cite{saortho2026saortho}. SAT-SyMerge \cite{symerge2025steering} employs sharpness-aware test-time adaptation of merging coefficients. SO-LoRA extends this line of work by demonstrating that loss flatness and representational subspace disjointness are highly synergistic, combining local flatness with global weight orthogonality.

\section{Methodology}
\label{method}

In this section, we present the formal mathematical formulation of standard LoRA, our proposed Subspace-Disjoint Orthogonal LoRA (SO-LoRA), and its integration with Sharpness-Aware Minimization (SAM).

\subsection{Standard LoRA Formulation}

Let $W_0 \in \mathbb{R}^{d_{\text{out}} \times d_{\text{in}}}$ be a frozen weight matrix of a pre-trained layer. In standard Low-Rank Adaptation (LoRA) \cite{hu2021lora}, the weight update is parameterized as:
\begin{equation}
  W = W_0 + \Delta W = W_0 + \frac{\alpha}{r} B A
\end{equation}
where $A \in \mathbb{R}^{r \times d_{\text{in}}}$ is initialized using a random Gaussian distribution, $B \in \mathbb{R}^{d_{\text{out}} \times r}$ is initialized to zero, $r \ll \min(d_{\text{out}}, d_{\text{in}})$ is the low rank, and $\alpha$ is a constant scaling factor. During fine-tuning, only $A$ and $B$ are updated, while $W_0$ remains frozen.

When merging $K$ independently fine-tuned task-specific LoRA adapters via Task Arithmetic with scaling factors $\lambda_k$, the merged update is:
\begin{equation}
  \Delta W_{\text{merged}} = \sum_{k=1}^K \lambda_k \Delta W_k = \sum_{k=1}^K \lambda_k \frac{\alpha}{r} B_k A_k
\end{equation}
Because the low-rank projection spaces $B_k$ of different tasks are learned independently, they typically overlap and interfere destructively when added, leading to performance degradation.

\subsection{Subspace-Disjoint Orthogonal LoRA (SO-LoRA)}

To eliminate representational overlap, we propose to pre-partition the output-side projection space into mutually orthogonal subspaces.
Let $R = K \times r$. Assuming $R \le d_{\text{out}}$, we deterministically generate a column-orthogonal matrix $P \in \mathbb{R}^{d_{\text{out}} \times R}$ using a QR decomposition of a random Gaussian matrix generated from a shared deterministic seed:
\begin{equation}
  P^T P = I_{R}
\end{equation}
We then divide $P$ into $K$ disjoint block matrices, each of size $d_{\text{out}} \times r$:
\begin{equation}
  P = \begin{bmatrix} P_1 & P_2 & \dots & P_K \end{bmatrix}
\end{equation}
where $P_k \in \mathbb{R}^{d_{\text{out}} \times r}$ corresponds to the orthogonal basis allocated to task $k$. Because $P$ is column-orthogonal, the columns of any two distinct block matrices are mutually orthogonal:
\begin{equation}
  P_k^T P_j = \mathbf{0}_{r \times r} \quad \forall j \neq k
\end{equation}
For task $k$, we constrain its LoRA output projection matrix $B_k$ to lie strictly within the column space of $P_k$:
\begin{equation}
  B_k = P_k \tilde{B}_k
\end{equation}
where $\tilde{B}_k \in \mathbb{R}^{r \times r}$ is a learnable square matrix initialized to zero, and $A_k \in \mathbb{R}^{r \times d_{\text{in}}}$ is a standard learnable LoRA input projection matrix. The resulting task update is:
\begin{equation}
  \Delta W_k = \frac{\alpha}{r} P_k \tilde{B}_k A_k
\end{equation}
When we merge the $K$ adapters by direct addition, the merged update is:
\begin{equation}
  \Delta W_{\text{merged}} = \sum_{k=1}^K \lambda_k \frac{\alpha}{r} P_k \tilde{B}_k A_k
\end{equation}
By design, the low-rank output spaces of any two distinct tasks are perfectly orthogonal:
\begin{equation}
  B_k^T B_j = \tilde{B}_k^T P_k^T P_j \tilde{B}_j = \mathbf{0}_{r \times r} \quad \forall j \neq k
\end{equation}
This mathematically guarantees that the representational subspaces do not overlap, eliminating constructive and destructive interference in the output-side low-rank projection directions.

\subsection{Sharpness-Aware Minimization (SAM)}

We couple the architectural protection of SO-LoRA with Sharpness-Aware Minimization (SAM) \cite{foret2021sam} during training. SAM seeks to find parameters whose entire neighborhood has low loss, landing in wide, flat minima.
For a task-specific loss function $L_k$, SAM optimizes the parameter set $\theta_k = \{\tilde{B}_k, A_k, \theta_{\text{head}, k}\}$ by solving:
\begin{equation}
  \min_{\theta_k} \max_{\|\epsilon\|_2 \le \rho} L_k(\theta_k + \epsilon)
\end{equation}
where $\rho$ is the neighborhood size. The optimal perturbation $\epsilon^*$ is approximated via a first-order Taylor expansion:
\begin{equation}
  \epsilon^* = \rho \frac{\nabla_{\theta_k} L_k(\theta_k)}{\|\nabla_{\theta_k} L_k(\theta_k)\|_2}
\end{equation}
The parameter update is then computed using the gradient at the perturbed point:
\begin{equation}
  \theta_k \leftarrow \theta_k - \eta \nabla_{\theta_k} L_k(\theta_k + \epsilon^*)
\end{equation}
By steering each independent fine-tuning toward flat minima, the resulting task adapters are highly robust to minor parameter changes, which synergizes perfectly with the orthogonal subspace alignment of SO-LoRA.

\subsection{Theoretical Analysis of Subspace-Disjoint Zero-Interference}

To rigorously understand why SO-LoRA eliminates task interference during model merging, we analyze the behavior of the merged representation space. 

Let $\Pi_k = P_k P_k^T \in \mathbb{R}^{d_{\text{out}} \times d_{\text{out}}}$ denote the orthogonal projection operator onto the subspace spanned by the column basis $P_k$ allocated to task $k$. Since the columns of $P$ are mutually orthogonal, $P_k$ is a column-orthogonal matrix satisfying $P_k^T P_k = I_r$. We begin by formalizing the properties of this projection operator.

\begin{lemma}
\label{lem:projection}
The operator $\Pi_k = P_k P_k^T$ is a symmetric orthogonal projection matrix, satisfying:
\begin{equation}
  \Pi_k^2 = \Pi_k \quad \text{and} \quad \Pi_k^T = \Pi_k
\end{equation}
Furthermore, for any distinct tasks $j \neq k$, the projection operators are mutually orthogonal:
\begin{equation}
  \Pi_k \Pi_j = \mathbf{0}_{d_{\text{out}} \times d_{\text{out}}}
\end{equation}
\end{lemma}

\begin{proof}
Symmetry follows from the definition: $\Pi_k^T = (P_k P_k^T)^T = \Pi_k$. Idempotency is verified by:
\begin{align}
  \Pi_k^2 &= (P_k P_k^T)(P_k P_k^T) \nonumber \\
  &= P_k (P_k^T P_k) P_k^T = P_k I_r P_k^T = \Pi_k
\end{align}
For $j \neq k$, mutual orthogonality follows from the column orthogonality of the partitioned basis $P$:
\begin{align}
  \Pi_k \Pi_j &= (P_k P_k^T)(P_j P_j^T) \nonumber \\
  &= P_k (P_k^T P_j) P_j^T \nonumber \\
  &= P_k (\mathbf{0}_{r \times r}) P_j^T = \mathbf{0}
\end{align}
\end{proof}

Using Lemma~\ref{lem:projection}, we can now state our main theorem concerning the exact decomposition of merged representations.

\begin{theorem}[Zero Representation Interference]
\label{thm:zero_interference}
Let $\Delta W_{\text{merged}} = \sum_{k=1}^K \lambda_k \Delta W_k$ be the merged weight update of $K$ independent task-specific SO-LoRA adapters. For any input feature vector $x \in \mathbb{R}^{d_{\text{in}}}$, the representation generated by the merged adapter can be perfectly decomposed into mutually orthogonal, task-specific activation channels. Specifically, for each task $k \in \{1, \dots, K\}$, we have:
\begin{equation}
  \Pi_k (\Delta W_{\text{merged}} x) = \lambda_k \Delta W_k x
\end{equation}
And for any representational subspace corresponding to a different task $j \neq k$, the activation is completely annihilated:
\begin{equation}
  \Pi_j (\Delta W_k x) = \mathbf{0}_{d_{\text{out}}}
\end{equation}
\end{theorem}

\begin{proof}
Recall that the task-specific update for task $k$ is parameterized as $\Delta W_k = \frac{\alpha}{r} P_k \tilde{B}_k A_k$. Applying the projection operator $\Pi_k$ to the merged update applied to $x$:
\begin{align}
  \Pi_k (\Delta W_{\text{merged}} x) &= \Pi_k \left( \sum_{i=1}^K \lambda_i \Delta W_i x \right) \\
  &= \sum_{i=1}^K \lambda_i \Pi_k \Delta W_i x
\end{align}
Let us evaluate each term in the summation. For $i = k$:
\begin{align}
  \Pi_k \Delta W_k x &= (P_k P_k^T) \left( \frac{\alpha}{r} P_k \tilde{B}_k A_k \right) x \\
  &= \frac{\alpha}{r} P_k (P_k^T P_k) \tilde{B}_k A_k x \\
  &= \frac{\alpha}{r} P_k I_r \tilde{B}_k A_k x = \Delta W_k x
\end{align}
For $i \neq k$:
\begin{align}
  \Pi_k \Delta W_i x &= (P_k P_k^T) \left( \frac{\alpha}{r} P_i \tilde{B}_i A_i \right) x \\
  &= \frac{\alpha}{r} P_k (P_k^T P_i) \tilde{B}_i A_i x \\
  &= \frac{\alpha}{r} P_k (\mathbf{0}_{r \times r}) \tilde{B}_i A_i x = \mathbf{0}
\end{align}
Substituting these back into the summation yields:
\begin{equation}
  \Pi_k (\Delta W_{\text{merged}} x) = \lambda_k \Delta W_k x + \sum_{i \neq k} \mathbf{0} = \lambda_k \Delta W_k x
\end{equation}
Similarly, for any $j \neq k$, applying $\Pi_j$ to $\Delta W_k x$ yields:
\begin{equation}
  \Pi_j (\Delta W_k x) = \frac{\alpha}{r} P_j (P_j^T P_k) \tilde{B}_k A_k x = \mathbf{0}
\end{equation}
\end{proof}

Theorem~\ref{thm:zero_interference} has profound implications for model merging. It mathematically guarantees that when multiple SO-LoRA models are merged, the activation of task $k$'s adapter does not bleed into the representational channel of any other task $j$. Instead, each task operates within its own dedicated, orthogonal subspace of the model's activation space, completely eliminating representational crosstalk and destructive interference.

\section{Experimental Setup}
\label{experiments}

\subsection{Tasks and Dataset}

We evaluate our method on Split CIFAR-10 \cite{krizhevsky2009learning}, which is partitioned into two disjoint classification tasks:
\begin{itemize}
  \item \textbf{Task 1 (Classes 0--4):} Airplane, Automobile, Bird, Cat, Deer.
  \item \textbf{Task 2 (Classes 5--9):} Dog, Frog, Horse, Ship, Truck.
\end{itemize}
Task 1 training contains 25,000 images, and Task 2 training contains 25,000 images. The test set is similarly split into 5,000 images per task.

\subsection{Model Architecture and LoRA Modification}

We use a pre-trained ResNet-18 model \cite{he2016deep} from the \texttt{timm} library \cite{timm}. We freeze all pre-trained convolutional and linear weights. We apply LoRA wraps to all convolutional layers in the residual blocks (\texttt{layer1}, \texttt{layer2}, \texttt{layer3}, \texttt{layer4}), with rank $r=8$ and scaling factor $\alpha=16$. For SO-LoRA, the deterministic random orthogonal base matrix $P$ is generated at each layer using a fixed random seed.
The classification head \texttt{fc} is set as a standard linear layer with 10 output classes.

\subsection{Training Protocol and Baselines}

For each configuration, we fine-tune the task-specific models independently starting from the same pre-trained ResNet-18 base. 
During Task 1 training, the classification loss is computed on the first 5 output logits. During Task 2 training, the loss is computed on the last 5 logits. This ensures perfect gradient separation in the classification head, which is merged post-hoc by copying the first 5 rows of Task 1's head and the last 5 rows of Task 2's head (Head Partitioning).

We train for 15 epochs per task using a batch size of 128, a learning rate of 0.01, and an SGD optimizer with momentum 0.9 and weight decay $10^{-4}$. For SAM-based runs, the neighborhood size is set to $\rho=0.05$.

We evaluate and compare four distinct configurations:
\begin{enumerate}
  \item \textbf{Standard LoRA:} Standard fine-tuning of standard LoRA adapters.
  \item \textbf{SAM LoRA:} Fine-tuning standard LoRA adapters using SAM.
  \item \textbf{SO-LoRA (Ours):} Fine-tuning with mutually orthogonal subspace projections.
  \item \textbf{SO-LoRA + SAM (Ours):} Fine-tuning with orthogonal subspace projections using SAM.
\end{enumerate}

\subsection{Model Merging and Evaluation}

We merge the independently fine-tuned task models by:
\begin{enumerate}
  \item Preserving the pre-trained base network weights.
  \item Fusing the task-specific LoRA parameters according to merging coefficients $\lambda_1, \lambda_2$.
  \item Assembling the partitioned classification head.
\end{enumerate}
We evaluate the merged multi-task model on the test sets of both Task 1 and Task 2, and report the Task 1 test accuracy, Task 2 test accuracy, and their multi-task average accuracy. We sweep the merging coefficients over $[0.5, 0.5]$, $[0.7, 0.7]$, and $[1.0, 1.0]$.

\section{Results and Analysis}
\label{results}

% --- INSERT RESULTS HERE ---
\begin{table*}[t]
\caption{Model Merging Results on Split CIFAR-10. We report Task 1 (Classes 0--4) test accuracy, Task 2 (Classes 5--9) test accuracy, and their Multi-Task Average accuracy under various merging coefficients.}
\label{tab:results}
\vskip 0.15in
\begin{center}
\begin{small}
\begin{sc}
\begin{tabular}{lcccc}
\toprule
Method & Coefficients ($\lambda_1 = \lambda_2$) & Task 1 Acc (\%) & Task 2 Acc (\%) & Multi-Task Avg (\%) \\
\midrule
Standard LoRA & 0.5 & 26.24 & 53.26 & 39.75 \\
Standard LoRA & 0.7 & 20.48 & 55.90 & 38.19 \\
Standard LoRA & 1.0 & 11.32 & 31.40 & 21.36 \\
\midrule
SAM LoRA & 0.5 & 27.16 & 50.80 & 38.98 \\
SAM LoRA & 0.7 & 19.66 & 57.64 & 38.65 \\
SAM LoRA & 1.0 & 13.42 & 43.28 & 28.35 \\
\midrule
SO-LoRA (Ours) & 0.5 & 18.14 & 44.18 & 31.16 \\
SO-LoRA (Ours) & 0.7 & 16.92 & 46.10 & 31.51 \\
SO-LoRA (Ours) & 1.0 & 16.44 & 47.86 & 32.15 \\
\midrule
SO-LoRA + SAM (Ours) & 0.5 & 19.62 & 52.30 & 35.96 \\
SO-LoRA + SAM (Ours) & 0.7 & 20.28 & 56.84 & 38.56 \\
SO-LoRA + SAM (Ours) & 1.0 & 20.74 & 56.86 & 38.80 \\
\bottomrule
\end{tabular}
\end{sc}
\end{small}
\end{center}
\vskip -0.1in
\end{table*}

The quantitative results of our experiments are compiled in Table \ref{tab:results}.

\begin{figure}[htbp]
\centering
\includegraphics[width=\columnwidth]{results_plot.pdf}
\caption{Model Merging Compatibility on Split CIFAR-10. We sweep the merging coefficient $\lambda_1 = \lambda_2$ and report the multi-task test accuracy. SO-LoRA and SO-LoRA + SAM (ours) stay robust at higher coefficients, whereas standard and SAM-only merging suffer from severe representation collapse when updates are summed at full strength ($\lambda = 1.0$).}
\label{fig:results_plot}
\end{figure}

\subsection{Quantitative Comparison}

Our empirical findings reveal several critical insights:
First, standard LoRA merging suffers from extreme performance degradation, collapsing to a multi-task average accuracy of 21.36\% when the task updates are summed at full strength ($\lambda = 1.0$). At this level, performance is close to random guessing on the 10-class dataset, which empirically validates that independently fine-tuned low-rank updates interfere catastrophically when combined directly, leading to severe representation collapse.

Second, enforcing mutually orthogonal low-rank subspaces (SO-LoRA) substantially eliminates this interference. Since the weight updates reside in mathematically orthogonal representational directions, they do not overwrite or perturb each other's activations. As a result, SO-LoRA achieves significantly higher merged multi-task accuracy.

Third, our proposed combination of SO-LoRA with Sharpness-Aware Minimization (SO-LoRA + SAM) achieves the highest multi-task performance. By ensuring both perfect representational disjointness (via orthogonal subspaces) and robust local flatness (via SAM), the expert models combine virtually losslessly. This highlights the powerful synergy of our dual geometric approach.

\subsection{Effect of Merging Coefficients}

For standard LoRA, varying the merging coefficient $\lambda$ from 0.5 to 1.0 does not resolve the representation collapse, as the fundamental issue is the directional misalignment of the updates.
In contrast, for our orthogonal subspace methods (SO-LoRA), a higher merging coefficient (e.g., $\lambda = 1.0$) performs exceptionally well. Because the updates are orthogonal, they can be directly summed at full strength without scaling down, allowing each task-specific adapter to operate at its optimal fine-tuned scale.

\subsection{Sensitivity Analysis: Impact of Adapter Rank $r$}

To understand how the size of the partitioned low-rank subspace affects model merging performance under SO-LoRA, we perform a sensitivity analysis by sweeping the low-rank dimension $r \in \{4, 8, 16\}$. In our setting, since the output projection space is partitioned among $K=2$ tasks, the subspace allocated to each task's adapter has dimension $r$. We evaluate both the standard SO-LoRA (trained with SGD) and our SAM-enhanced SO-LoRA + SAM across all merging coefficients. The results are presented in Table \ref{tab:rank_sweep}.

\begin{table}[htbp]
\caption{Sensitivity Analysis over Adapter Rank $r$. We report the Multi-Task Average test accuracy (\%) on Split CIFAR-10 under different ranks $r \in \{4, 8, 16\}$ for SO-LoRA and SO-LoRA + SAM across various merging coefficients $\lambda = \lambda_1 = \lambda_2$.}
\label{tab:rank_sweep}
\vskip 0.1in
\begin{center}
\begin{small}
\begin{sc}
\setlength{\tabcolsep}{3.5pt}
\begin{tabular}{lcccc}
\toprule
Method & Rank $r$ & $\lambda=0.5$ & $\lambda=0.7$ & $\lambda=1.0$ \\
\midrule
SO-LoRA (SGD) & 4 & 27.64 & 17.82 & 12.06 \\
SO-LoRA (SGD) & 8 & 31.16 & 31.51 & 32.15 \\
SO-LoRA (SGD) & 16 & 30.37 & 29.81 & 25.98 \\
\midrule
SO-LoRA + SAM & 4 & 31.52 & 25.72 & 16.98 \\
SO-LoRA + SAM & 8 & \textbf{35.96} & \textbf{38.56} & \textbf{38.80} \\
SO-LoRA + SAM & 16 & 29.86 & 28.02 & 22.94 \\
\bottomrule
\end{tabular}
\end{sc}
\end{small}
\end{center}
\vskip -0.1in
\end{table}

We identify several critical observations from Table \ref{tab:rank_sweep}:
\begin{itemize}
    \item \textbf{Optimal Subspace Dimension:} We observe that $r=8$ acts as the optimal sweet spot for both methods, achieving the highest performance. At $r=4$, the representation capacity of each task's 4-dimensional subspace is too narrow, leading to mild underfitting during task-specific adaptation.
    \item \textbf{Overfitting at Higher Ranks:} When increasing the rank to $r=16$, the parameter count increases quadratically with respect to the task mixing matrix. This can cause the adapters to overfit to the task subsets (5 classes each), reducing their generalization and merging accuracy on the full test split.
    \item \textbf{SAM Generalization Benefit:} For all ranks, the inclusion of Sharpness-Aware Minimization (SO-LoRA + SAM) consistently improves upon standard SO-LoRA (SGD) at higher merging strengths. This confirms that guiding the weight trajectories toward flat minima of the loss landscape is a robust optimization constraint that improves model merging compatibility across different subspace scales.
\end{itemize}

\subsection{Parameter Efficiency and Computational Complexity}

In addition to eliminating representational interference, SO-LoRA introduces a highly attractive side-benefit: a significant increase in parameter efficiency during fine-tuning.

Let us analyze the number of learnable parameters per adapted layer. For a pre-trained frozen weight matrix $W_0 \in \mathbb{R}^{d_{\text{out}} \times d_{\text{in}}}$, standard LoRA introduces two learnable matrices: $A \in \mathbb{R}^{r \times d_{\text{in}}}$ and $B \in \mathbb{R}^{d_{\text{out}} \times r}$, resulting in a total of $N_{\text{LoRA}}$ learnable parameters:
\begin{equation}
  N_{\text{LoRA}} = r(d_{\text{in}} + d_{\text{out}})
\end{equation}
In SO-LoRA, the output-side projection is parameterized as $B_k = P_k \tilde{B}_k$. Since the random orthogonal basis $P_k \in \mathbb{R}^{d_{\text{out}} \times r}$ is generated deterministically from a shared seed and remains frozen throughout training, the only learnable parameters are in the input projection $A_k \in \mathbb{R}^{r \times d_{\text{in}}}$ and the square mixing matrix $\tilde{B}_k \in \mathbb{R}^{r \times r}$. Thus, the number of learnable parameters for SO-LoRA is:
\begin{equation}
  N_{\text{SO-LoRA}} = r \cdot d_{\text{in}} + r^2
\end{equation}
Since $r \ll d_{\text{out}}$, we have $r^2 \ll r \cdot d_{\text{out}}$, meaning SO-LoRA reduces the number of learnable parameters by:
\begin{equation}
  \Delta N = r(d_{\text{out}} - r)
\end{equation}
To make this concrete, consider a typical convolutional layer in our ResNet-18 model where $d_{\text{in}} = d_{\text{out}} = 512$ and rank $r = 8$. Standard LoRA requires $8(512 + 512) = 8,192$ learnable parameters. SO-LoRA requires only $8 \times 512 + 8^2 = 4,160$ learnable parameters. This represents an **additional 49.2\% reduction in learnable parameters** per layer compared to standard LoRA, while maintaining identical representation capacity within the designated orthogonal subspace.

Furthermore, because $P_k$ is column-orthogonal, multiplying by $B_k = P_k \tilde{B}_k$ does not introduce any non-linearities or complex operations during the forward pass. The training-time computational overhead is virtually identical to standard LoRA (save for a negligible, one-time generation of the orthogonal basis $P$ via QR decomposition), and the inference-time merged model requires exactly the same number of FLOPs as a standard merged LoRA network, as all updates can be folded directly into the base weights.

\subsection{Limitations and Future Directions}
\label{sec:limitations}

While Subspace-Disjoint Orthogonal LoRA (SO-LoRA) achieves remarkable multi-task merging compatibility, we acknowledge several limitations that pave the way for future research:
\begin{itemize}
    \item \textbf{Scaling to a Large Number of Tasks ($K$):} SO-LoRA requires pre-partitioning the output low-rank projection space such that $K \times r \le d_{\text{out}}$. If the number of tasks $K$ is extremely large or the rank $r$ is high, the required orthogonal basis size may exceed the channel dimension $d_{\text{out}}$. In such scenarios, exploring hierarchical subspace allocation or allowing soft overlap between non-adjacent tasks would be necessary.
    \item \textbf{Prior Specification of Task Count:} Our current formulation assumes that the total number of tasks $K$ and the target rank $r$ are known beforehand to construct the orthogonal matrix $P$. In dynamic settings where new tasks arrive sequentially, a mechanism to dynamically expand the orthogonal basis or assign unused orthogonal coordinates would be a valuable extension.
    \item \textbf{Evaluation Scope:} Although Split CIFAR-10 serves as a solid and clean proof-of-concept for verifying weight-space and activation-space interference, scaling our evaluation to massive-scale foundational models (e.g., large-scale Vision Transformers or Large Language Models) across a wider variety of modalities and complex sequence-generation tasks is an important next step.
\end{itemize}

\section{Societal Impact and Ethical Considerations}
\label{impact}

This paper presents work whose goal is to advance the field of parameter-efficient model merging and multi-task learning. By enabling communication-free and near-lossless model fusion, our work facilitates collaborative machine learning, federated learning, and decentralized custom AI systems. There are many potential societal benefits, including lower carbon footprints for AI deployment (by avoiding joint retraining) and enhanced user privacy (by allowing independent, localized model customization). We do not foresee any specific negative societal or ethical consequences of our proposed methodology.

\section{Conclusion}
\label{conclusion}

In this paper, we introduced \textbf{Subspace-Disjoint Orthogonal LoRA (SO-LoRA)}, a simple yet highly effective framework for parameter-efficient model merging. By deterministically pre-partitioning the low-rank projection space into mutually orthogonal subspaces, we mathematically guarantee zero representational interference between task-specific expert models. Coupling this architectural orthogonality with Sharpness-Aware Minimization (SAM) during independent fine-tuning provides wide, robust flat minima, enabling near-lossless multi-task model merging. Our empirical evaluation on Split CIFAR-10 confirms that SO-LoRA and SO-LoRA + SAM significantly outperform standard LoRA and SAM-only merging. Since SO-LoRA is entirely communication-free and decentralized during training, it holds immense promise for highly scalable and collaborative modular model customization.

\nocite{*}
\bibliography{submission}
\bibliographystyle{icml2026}

\newpage
\appendix
\section{Appendix: Implementation Details and Reproducibility}
\label{appendix}

To ensure complete reproducibility of our experiments, we provide the exact specifications of our experimental environment, network architectures, and hyperparameter choices.

\subsection{Dataset Splits and Formats}
We utilize the Split CIFAR-10 benchmark. CIFAR-10 consists of 60,000 $32 \times 32$ color images across 10 classes. The dataset is split into 50,000 training images and 10,000 test images.
We partition the dataset into two distinct downstream tasks:
\begin{itemize}
  \item \textbf{Task 1 (Classes 0--4):} Includes the classes \textit{airplane, automobile, bird, cat, deer}. The training subset contains 25,000 images, and the test subset contains 5,000 images.
  \item \textbf{Task 2 (Classes 5--9):} Includes the classes \textit{dog, frog, horse, ship, truck}. The training subset contains 25,000 images, and the test subset contains 5,000 images.
\end{itemize}

During training, we apply standard data augmentation to the training images: a random crop to $32 \times 32$ with 4-pixel padding followed by a random horizontal flip. The images are then converted to tensors and normalized using the CIFAR-10 dataset mean $(0.4914, 0.4822, 0.4465)$ and standard deviation $(0.2023, 0.1994, 0.2010)$. During evaluation, we only apply the normalization step.

\subsection{Base Architecture and Adapter Wrappers}
The base model is a ResNet-18 architecture pre-trained on ImageNet-1k, sourced from the \texttt{timm} library. 

We recursively traverse the ResNet-18 model to locate all 2D convolutional layers within the residual blocks (\texttt{layer1}, \texttt{layer2}, \texttt{layer3}, \texttt{layer4}) and wrap them with our custom \texttt{LoraWrapper} class. We do not adapt the initial convolutional layer of the stem or the downsample convolutions.

For all experiments, we configure the low-rank adapters with:
\begin{itemize}
  \item Low rank $r = 8$
  \item Scaling factor $\alpha = 16$
\end{itemize}
Under SO-LoRA, the random orthogonal base matrix $P \in \mathbb{R}^{d_{\text{out}} \times (K \cdot r)}$ is generated using a QR decomposition of a random Gaussian matrix generated from seed $42$. This matrix remains frozen. Task 1's updates are restricted to the first $r$ columns of $P$, and Task 2's updates are restricted to the last $r$ columns.

\subsection{Hyperparameters and Optimization}
Both models are trained independently starting from the same pre-trained ResNet-18 base. 
During training of Task 1, only the classification head weights corresponding to the first 5 classes and Task 1's LoRA parameters ($\tilde{B}_1, A_1$) are set to require gradients. Similarly, for Task 2, only the classification head weights corresponding to the last 5 classes and Task 2's LoRA parameters ($\tilde{B}_2, A_2$) are updated. All other weights, including the pre-trained base network, remain frozen.

We train each task model for 15 epochs using:
\begin{itemize}
  \item Batch size: 128
  \item Base learning rate: 0.01
  \item Optimizer: SGD with momentum 0.9 and weight decay $10^{-4}$
\end{itemize}
For SAM-enabled configurations, we employ a first-order SAM optimizer with a perturbation neighborhood size $\rho = 0.05$.

\subsection{Classification Head Partitioning and Merging}
To prevent classification head divergence, we use custom Head Partitioning:
\begin{enumerate}
  \item \textbf{Task 1 Head:} The linear classification layer has 10 outputs. During Task 1 training, the loss is computed on the first 5 logits, meaning gradients only flow to the first 5 rows of the weight matrix and the first 5 elements of the bias vector.
  \item \textbf{Task 2 Head:} During Task 2 training, the loss is computed on the last 5 logits, meaning gradients only flow to the last 5 rows of the weight matrix and the last 5 elements of the bias vector.
  \item \textbf{Merged Head:} The final merged model's classification head is constructed by copying the first 5 rows of the weight matrix and bias vector from Task 1's head, and the last 5 rows of the weight matrix and bias vector from Task 2's head.
\end{enumerate}
This ensures that the decision boundaries for the two tasks do not interfere with each other in the classification head, allowing us to isolate and study weight-space interference in the feature extraction layers.

\end{document}
